\section{Conclusion and Future Work}

% klasik luenberger bu sistemlerde de durum kestiriyor
% disturbance için ek bilgi ışığında gözleyici tasarlanabilir 
% PFMnin ek bir disturbance modeli olmadan gauge invariance'a dayanan bir ayrıştırma yöntemi oldunu

% PFM output redundant sistemlerde tanımlı
% PFM esitmation olmadan linear mapping ayrıştırdığını
% PFM dinamik değil statik bir yapıda olduğunu
% PFM ile gürültü varken de ayrıştırma mümkün olduğunu 
% Eğer sistem Output redundant ise PFM klasik gözleyiciye bir alternatif olduğu gösterildi

% For estimation of states of output redundant systems, the Luenberger observer is a classical tool. This extends to disturbance estimation, if a model or additional assumptions are made for the disturbance. However, it has been shown that, for output redundant systems, a projection filtering method based on the gauge-invariance principle can be used to both reconstruct the states and disturbances. It has been noted that the separation is a static process, rather than an estimation process. It has been observed that, the static nature proves to be an important advantage, since this can be achieved with high gain Luenberger observers, but with a cost of higher noise amplification. It is also shown that, the PFM can be applied for outputs where noise is present. It is observed that the PFM is an alternative to classical observers for output redundant systems.

For estimation of states of output redundant systems, the Luenberger observer is a classical tool. This extends to disturbance estimation, if a model or additional assumptions are made for the disturbance. 

However, it has been shown that, for output redundant systems, a projection filtering method can be used to both reconstruct the states and disturbances. The PFM is based on the gauge-invariance principle, where the invariance can be chosen to isolate components for recovery under appropriate transformations, which then in turn is used to separate the linear mapping in the output. It has been noted that the separation is a static process, rather than an estimation process. It has been observed that, the static nature proves to be an important advantage, since this can be achieved with high gain Luenberger observers, but with a cost of higher noise amplification. 

It is also shown that, the PFM can be applied for outputs where noise is present. It is observed that the PFM is an alternative to classical observers for output redundant systems.

A natural extension of this work is to apply PFM to systems with nonlinear dynamics and output mappings, etc. 

